% concatenated from main.tex
\documentclass[12pt,a4paper,reqno]{amsart}
\usepackage[utf8]{inputenc}
\usepackage[english]{babel}
%\usepackage{amsrefs}
\usepackage{enumerate}
%\usepackage{refcheck, hyperref}
\usepackage{mathtools}
\usepackage{mathrsfs}
\usepackage{soul, xcolor}
\usepackage{wasysym}
\usepackage{amsmath,amssymb,xpatch,relsize,graphics,graphicx,amsthm}
%\usepackage[square,numbers]{natbib}
%\bibliographystyle{abbrvnat}
%\usepackage{acta-m,euler}
\usepackage{float,graphicx}
\usepackage{subcaption}

%\usepackage[demo]{graphicx}
%\usepackage{subfig}
%\usepackage{hyperref}
%\usepackage{refcheck}
%\usepackage{soul}
\DeclareMathOperator{\RE}{Re}
\DeclareMathOperator{\IM}{Im}


\xapptocmd\normalsize{%
	\abovedisplayskip=12pt plus 3pt minus 9pt
	\abovedisplayshortskip=0pt plus 3pt
	\belowdisplayskip=12pt plus 3pt minus 9pt
	\belowdisplayshortskip=7pt plus 3pt minus 4pt
}{}{}
\textwidth=465pt \evensidemargin=-15pt \oddsidemargin=0pt
\marginparsep=8pt \marginparpush=8pt \textheight=660pt
\topmargin=-20pt

\theoremstyle{definition}
\newtheorem{definition}{Definition}[section]
\newtheorem{remark}[definition]{Remark}
\newtheorem{example}[definition]{Example}
\theoremstyle{plain}
\newtheorem{theorem}[definition]{Theorem}
\newtheorem{corollary}[definition]{Corollary}
\newtheorem{lemma}[definition]{Lemma}
\newtheorem{proposition}[definition]{Proposition}

\numberwithin{equation}{section}

\title[Sharp Bohr-Type  inequalities]{Sharp Bohr-Type  inequalities for certain classes of close-to-convex  functions}
\author[S. Rana]{Shalini Rana}
\address{Department of Mathematics, University of Delhi, Delhi--110 007, India}
\email{shalinirana3010@gmail.com}
\author[N.K. Jain]{Naveen Kumar Jain*}
\address {Department of Mathematics, Aryabhatta College, Delhi-110021,India}
\email{naveenjain@aryabhattacollege.ac.in{*}}

\date{}

\keywords{Starlike functions;  Analytic Functions;  Bohr Inequality; Subordination\\}

\subjclass[2010]{30C45, 30C50, 30C80}

%$*$-Corresponding Author
\allowdisplaybreaks

\begin{document}
%\author[add1]{Shalini Rana}
%\ead{shalinirana3010@gmail.com}
%\author[add3]{Naveen Kumar Jain}
%\ead{naveenjain05@gmail.com}

%\address[add1]{Department of Mathematics, University of Delhi, Delhi 110007, India}
%\address[add3]{Department of Mathematics, Aryabhatta College, Delhi 110021, India}
\begin{abstract}
In this article, we determine the Rogosinski radii for certain subclasses of close-to-convex functions defined on open unit disc $\mathbb{D}= \{z \in \mathbb{C}: |z| < 1\}$. Furthermore, we establish improved versions of the classical Bohr inequality and the Bohr-Rogosinski inequality pertaining to these subclasses. We demonstrate that all results derived in the study are sharp.
\end{abstract}
\maketitle




\section{Introduction and preliminaries}
Let $\mathcal{A}$ be the class of analytic functions $g$ with the Taylor series expansion $g(z)=\sum_{n=0}^{\infty}a_nz^n$ in the unit disc $\mathbb{D}=\{ z \in \mathbb{C}: |z|<1 \}$. The classical Bohr inequality \cite{boh1914} says that if $g\in\mathcal{A}$  and  $g$ maps the unit disc into itself, then
\begin{equation}\label{first}
\sum_{n=0}^{\infty}|a_n|r^n\leq1\quad\text{for}\quad |z|=r\leq\frac{1}{3}
\end{equation}
and the radius  $1/3$ cannot be improved. The radius $1/3$ is known as Bohr radius, while the inequality \eqref{first} is known as Bohr inequality. Initially, the inequality \eqref{first} was established by Harald Bohr for radius $r\leq1/6$ and later Weiner, Riesz and Schur obtained the inequality for $r\leq1/3$. It is easy to see that $(1-|a_0|)$ is equal to $d(g(0), \partial g(\mathbb{D}))$, where $d$ is the Euclidean distance and $\partial g(\mathbb{D})$  is the boundary of $g(\mathbb{D})$. Therefore, the Bohr inequality can be re-written in the following Euclidean distance form
\begin{equation}\label{second}
d\left(\sum_{n=0}^{\infty}|a_nz^n|, |a_0|\right)=\sum_{n=1}^{\infty}|a_nz^n|\leq1-|a_0|=d(g(0), \partial g(\mathbb{D})).
\end{equation}
The inequality \eqref{second} can be generalized to class $\mathcal{G}$ consisting of functions $g$ analytic in $\mathbb{D}$ such that $g(\mathbb{D})\subset\Delta$, where $\Delta$ is any given domain. The class $\mathcal{G}$ satisfies the Bohr phenomenon if there exists largest $r_{\Delta}\in (0, 1)$ such that
\begin{equation*}%\label{eq14}
d\left(\sum_{n=0}^{\infty}|a_nz^n|, |a_0|\right)=\sum_{n=1}^{\infty}|a_nz^n|\leq d(g(0), \partial \Delta)
\end{equation*}
holds for all $|z|=r\leq r_{\Delta}$ and for all the functions $g\in\mathcal{G}$. It is interestingly
enough, for the convex domain $\Delta$, the radius $r_{\Delta}=1/3$ satisfying \eqref{second} coincides with Bohr radius  \cite{aiz2007}. The Bohr radius has been obtained for the class $\mathcal{G}$ when $\Delta$ has positive real part \cite{Ali2019} and  Janowski starlike \cite{Ana2021} (see also \cite{Ahu2020}). Several  Bohr radii have been obtained for the classes of analytic functions $g$ satisfying the differential
subordination $g(z)+\beta zg'(z)+\gamma z^2g''(z)\prec k(z)$, when the function $k$ is convex and starlike \cite{Abu2014}, convex and starlike of order $\alpha$ \cite{Jai2020} and  Janowski starlike \cite{Ana2021}.
\\
\noindent
The Bohr inequality has emerged as an active research area for operator algebraists after Dixon \cite{Dix1995} established a connection between the inequality and the characterization of Banach Algebra that satisfies Von Neumann inequality. Improved versions of Bohr inequality for various classes of analytic functions  becomes now a days an active research area. Kayumov et al. \cite{Kay2018} obtained several different improved versions of the classical Bohr inequality. Recently, Ahamed and  Ahammed \cite{ahm2023} obtain sharp refined Bohr–Rogosinski-type
inequalities for the subordination class of univalent and convex function. In 2021, Liu et al. \cite{LIU2021} generalized and improved several Bohr-type inequalities for bounded analytic functions. Recently, Rana and Jain  \cite{sha2025} has established several improved Bohr-type inequalities for starlike functions of order $\alpha$,  functions whose derivative has positive real part and the  functions  convex in the direction of imaginary axis, see also \cite{Pon2020, Hua2020, Liu2023, Pon2023}.
\\
\noindent
In 1952, Kaplan \cite{Kap1952} introduced the class of close-to-convex functions $f\in\mathcal{A}$ satisfying the condition
\begin{equation}\label{eq23}
Re \left(\frac{zf^{'}}{g} \right) >0\quad (z \in \mathbb{D}),
\end{equation}
where $g$ is a function starlike  with respect to origin  of the form
\begin{equation}\label{eq24}
g(z)=b_1z+...  Re(b_1)> 0.
\end{equation}
The class is denoted by $\mathcal{C}$. Later, Silverman and Telage \cite{sil1979} introduced the three subclasses, $\mathcal{C}_1$, $\mathcal{C}_2$ and $\mathcal{C}_3$ of class $\mathcal{C}$.
\noindent
A function $f$ is said to be in the class $\mathcal{C}_1$ if there exists a convex function $g$ of the form \eqref{eq24} such that \eqref{eq23} is satisfied.
A function $f$ is said to be in the class $\mathcal{C}_2$ if there exists a convex function $g$ of the form \eqref{eq24} satisfying
\begin{equation*}
Re\left( \frac{(zf^{'})^{'}}{g^{'}}  \right)>0\quad (z \in \mathbb{D}).
\end{equation*}
\noindent
A function $f$ is said to be in the class $\mathcal{C}_3$ if
\begin{equation*}
Re\left(   \frac{z((zf^{'})^{'})^{'}}{(zg^{'})^{'}}  \right)>0 \quad (z \in \mathbb{D}),
\end{equation*}
where  $g$ is a convex function of the form \eqref{eq24}. They proved that  $\mathcal{C}_3 \subset \mathcal{C}_2\subset \mathcal{C}_1 \subset \mathcal{C}$.
In 2017,
Kayumov and Ponnusamy \cite{kay2017} introduced and found the following Bohr-Rogosinski inequality and  Bohr–Rogosinski radius for the class of  functions mapping unit disc to unit disc.
\begin{theorem}\cite{kay2017}
Suppose that $f:\mathbb{D}\rightarrow\mathbb{D}$ and $f(z)=\sum_{n=0}^{\infty}a_nz^n$. Then
\begin{equation*}
|f(z)|+\sum_{n=N}^{\infty}|a_n r^n|\leq1 \quad \text{for}\, |z|=r\leq R_N,
\end{equation*}
where $R_N$ is the positive root of the equation $2(1+r)r^N-(1-r)^2=0$. The radius $R_N$ is the best possible. Moreover,
\begin{equation*}
|f(z)|^2+\sum_{n=N}^{\infty}|a_n r^n|\leq1 \quad \text{for} \,|z|=r\leq R'_N,
\end{equation*}
where $R'_N$ is the positive root of the equation $(1+r)r^N-(1-r)^2=0$. The radius $R'_N$ is the best possible.
\end{theorem}
\noindent
For the results on Bohr-Rogosinski inequality, we refer to the articles \cite{ALK2020, KAV2021, PON2022,Pon2020}.
\\
\noindent
In the present work, we extend this line of research by determining the sharp Bohr–Rogosinski radius for the classes $\mathcal{C}_1, \mathcal{C}_2$ and $\mathcal{C}_3$ together with  improved
versions of the classical Bohr’s inequality . In section 2, section 3, and section 4, we find the sharp Bohr–Rogosinski radius and its various versions for the classes $\mathcal{C}_1, \mathcal{C}_2$ and $\mathcal{C}_3$ respectively.


\begin{section}{Bohr-type inequalities for the class $\mathcal{C}_1$}
\noindent
We need the following lemmas for main results.
\begin{lemma}\label{lm1}\cite{Sil1972}
If $f(z)= z + \sum_{n=2}^{\infty} a_n z^n \in \mathcal{C}_1$, then
\begin{equation} \label{coef1}
|a_n| \le 2- \frac{1}{n}
\end{equation}
with equality for $k(z,1,-1)$, where
\begin{equation*}
k(z,x,y)= (1+x) \frac{z}{(1-yz)} + x \bar{y} \log{(1-yz)}
\end{equation*}
for $z \in \Delta$ and $|x|=|y|=1$.
\end{lemma}
\begin{lemma}\label{lm2}{\cite{Sil1972}}
If $f \in \mathcal{C}_1$ then
\begin{equation} \label{growth1}
\frac{2r}{1+r} -\log(1+r) \le |f(z)| \le \frac{2r}{1-r} +\log(1-r)
\end{equation}
and
\begin{equation} \label{distortion1}
\frac{1-r}{(1+r)^2} \le |f^{'}(z)| \le \frac{1+r}{(1-r)^2}
\end{equation}
for all $|z| \le r$ and equality holds for $k(z,1,1)$ at $z=\pm r$  where
\begin{equation*}
k(z,x,y)= (1+x) \frac{z}{(1-yz)} + x \bar{y} \log{(1-yz)}
\end{equation*}
for $z \in \Delta$ and $|x|=|y|=1$.
\end{lemma}
\noindent
We begin by determining the sharp Bohr radius associated with the Bohr inequality for analytic functions $f(z)$. In this modified form of the inequality, the classical coefficients $|a_0|$ and $|a_1|$ are replaced by $|f(z)|$ and $|f^{'}(z)|$, respectively, allowing us to explore a functional analogue of the traditional Bohr-type estimate.
\begin{theorem}
Suppose that $f(z)=z + \sum_{n=2}^{\infty} a_n z^n \in \mathcal{C}_1$, then
\begin{equation}\label{2.1.1}
|f(z)| + |f^{'}(z)||z| + \sum_{n=2}^{\infty} |a_n z^n| \le d( f(0), \partial{f(\mathbb{D})})
\end{equation}
for $|z| \le r_{11}$, where $r_{11}=0.110377$ is the solution of
\begin{equation*}
    \log2-1-r (-6+r (2+r-\log2)+\log4)+2 (1-r)^2 \log(1-r) =0.
\end{equation*}
The result is sharp.
\end{theorem}
\begin{proof}
We begin the proof by evaluating $d(f(0), \partial{f(\Delta)})$.
From growth estimate of $f(z) \in \mathcal{C}_1$ in equation ($\ref{growth1}$), we get
\begin{equation*}
d(f(0), \partial{f(\mathbb{D})}) = \liminf \limits_{|z| \rightarrow 1} |f(z)-f(0)| \ge \lim \limits_{r \rightarrow 1} \left(\frac{2r}{1+r}- \log(1+r)\right)
\end{equation*}
which gives
\begin{equation}\label{eq1}
d(f(0), \partial{f(\mathbb{D})}) \ge (1-\log{2}).
\end{equation}
For $|z|\leq r$, by $(\ref{coef1})$, $(\ref{growth1})$, $(\ref{distortion1})$ , we have
\begin{equation}\label{eq3}
|f(z)| + |f^{'}(z)||z| + \sum_{n=2}^{\infty} |a_n| |z|^n\leq \frac{2r}{1-r} + \log(1-r) + \frac{r(1+r)}{(1-r)^2} + \sum_{n=2}^{\infty} \left( 2-\frac{1}{n} \right) r^{n}
\end{equation}
Define a function $P:[0, 1)\rightarrow \mathbb{R}$ by
\begin{equation*}
P(r)=  \log2-1-r (-6+r (2+r-\log2)+\log4)+2 (1-r)^2 \log(1-r).
\end{equation*}
It can be easily seen that $P(0)=\log2-1 \simeq-0.306853< 0$ and $P(1/2)={11}/{8}-{\log2}/{4}\simeq1.20171 > 0$. Therefore by intermediate value theorem $P(r)$ has  a  root in $ (0,1/2)$. Next to show  the uniqueness of root in $(0, 1/2)$, it is sufficient to show that the function $P$ is a strictly increasing function in $(0,1/2)$.
 It is easy to see that  $P'(r)=(2+6 r-5 r^2+r^3)/{(1-r)^3}>(2+6 r-5 r+r^3)/{(1-r)^3}=(2+ r+r^3)/{(1-r)^3}>0$ for all $r \in (0,1)$, which shows that $P$ is strictly increasing function.  Therefore, $P(r)$ has the unique root in $(0, 1/2)$, say $r_{11}=0.110377$.
  Thus,
  \begin{equation}\label{eq4}
  P(r_{11})=  \log2-1-r_{11} (-6+r_{11} (2+r_{11}-\log2)+\log4)+2 (1-r_{11})^2 \log(1-r_{11})=0
  \end{equation}
   and $P(r)\leq0$ for $r\in(0, r_{11})$.
\noindent
 In view of the above argument together with   \eqref{eq1}   and \eqref{eq3}, for  $|z| \le r_{11}$, we have
\begin{align*}\label{eq2}
|f(z)| + &|f^{'}(z)||z| + \sum_{n=2}^{\infty} |a_n| |z|^n-d(f(0), \partial{f(\mathbb{D})})\\
& \le \frac{2r}{1-r} + \log(1-r) + \frac{r(1+r)}{(1-r)^2} + \sum_{n=2}^{\infty} \left( 2-\frac{1}{n} \right) r^{n}-(1-\log2)\\
&= \frac{\log2-1-r (-6+r (2+r-\log2)+\log4)+2 (1-r)^2 \log(1-r)}{(1 - r)^2}\\
&\leq0
\end{align*}
\begin{figure}
  \centering
  % Requires \usepackage{graphicx}
  \includegraphics[width=8cm]{Thm23.pdf}\\
  \caption{The radius $r_{11}=0.110377$ is sharp}%\label{}
\end{figure}
\noindent
To prove the sharpness, we assume
\begin{equation*}
    f(z)=\frac{2z}{(1+z)}- \log(1+z).
\end{equation*}
Note that the power series expansion of $f(z)$ is given by
\begin{equation*}
    f(z)= \sum_{n=1}^{\infty} (-1)^{n+1} \left(2-\frac{1}{n} \right) z^n.
\end{equation*}
At $z=-r=-r_{11}$, by \eqref{eq4}, we have
\begin{align*}
|f(z)| + |f^{'}(z)||z| + &\sum_{n=2}^{\infty} |a_n| |z|^n\\
&= \left|\frac{2z}{(1+z)}- \log(1+z) \right| +  \left| \frac{(1-z)z}{(1+z)^2} \right|  +   \sum_{n=2}^{\infty} \left| (-1)^{n+1} \left(2-\frac{1}{n} \right)z^n \right| \\
& =  \left|\frac{2(-r)}{(1-r)}- \log(1-r) \right| + \left| \frac{-r(1+r)}{(1-r)^2} \right|  +   \sum_{n=2}^{\infty}  \left|2-\frac{1}{n} \right| |-r|^n\\
&=  \frac{2r}{(1-r)}+\log(1-r) +  \frac{r(1+r)}{(1-r)^2}   +   \sum_{n=2}^{\infty}  \left(2-\frac{1}{n} \right) |r|^n\\
&=1-\log 2.\\
&=\liminf \limits_{|z| \rightarrow 1} \left(\frac{2z}{1+z}- \log(1+z)\right)\\
&=d(f(0), \partial{f(\mathbb{D})}).
\end{align*}
Hence, the result is sharp.
\end{proof}
\noindent
Next, by considering powers of the coefficient $(|a_n|)$, we evaluate the corresponding sharp Bohr radii for the Bohr inequality for class $\mathcal{C}_1$. This extension allows us to analyze how the powers of the moduli of the coefficients influence the Bohr-type bounds, thereby enriching the understanding of Bohr phenomena in the class $\mathcal{C}_1$.
\begin{theorem}
Suppose that $f(z)=z + \sum_{n=2}^{\infty} a_n z^n \in \mathcal{C}_1$, then
\begin{equation} \label{2.1.2}
|z|+ \sum_{n=2}^{\infty} |a_n z^n| + \sum_{n=2}^{\infty} |a_n|^{p} |z|^{np}  \le d(f(0), \partial{f(\mathbb{D})})
\end{equation}
for $|z| \le r_p$ and $p \ge 1$, where $r_p$ is the solution of
\begin{equation*}
(1-r)\sum_{n=2}^{\infty} \left( 2-\frac{1}{n} \right)^{p} r^{pn}- (1-3 r+r \log2-\log(2-2 r)+r \log(1-r))=0.
\end{equation*}
The result is sharp.
\end{theorem}
\begin{proof}
By Lemma \ref{lm2}, it follows that
\begin{equation}\label{eq5}
d(f(0), \partial{f(\mathbb{D})})= \liminf \limits_{|z| \rightarrow 1} |f(z)-f(0)|
 \ge \lim \limits_{r \rightarrow 1} \left(\frac{2r}{1+r}- \log(1+r)\right)
\end{equation}
which gives $d(f(0), \partial{f(\mathbb{D})}) \ge (1-\log{2})$ for $f \in \mathcal{C}_1$. By Lemma \ref{lm1},  for $|z| \le r$ and $p \ge 1$, we have
\begin{align*}
|z|+ \sum_{n=2}^{\infty} |a_n| |z|^n&+ \sum_{n=2}^{\infty} |a_n|^{p} |z|^{np} \le r+ \sum_{n=2}^{\infty} \left( 2-\frac{1}{n} \right) r^{n}+ \sum_{n=2}^{\infty} \left( 2-\frac{1}{n} \right)^{p} r^{pn}\\
&=\sum_{n=2}^{\infty} \left( 2-\frac{1}{n} \right)^{p} r^{pn}- \frac{(1-3 r+r \log2-\log(2-2 r)+r \log(1-r))}{1-r}.
\end{align*}
Let  $Q$ be a function  defined by
\begin{align*}
Q(r)=\sum_{n=2}^{\infty} \left( 2-\frac{1}{n} \right)^{p} r^{pn}- \frac{1-3 r+r \log2-\log(2-2 r)+r \log(1-r)}{1-r}, \quad r\in [0, 1).
\end{align*}
Since, $Q(0)=-(1-\log 2)<0$ and
 \[
 Q(1/2)
 = 1+ \sum_{n=2}^{\infty} \left( 2-\frac{1}{n} \right)^{p} \frac{1}{2^{pn}}> 0
\]
by the Intermediate Value Theorem it follows that $Q(r)=0$ has a root in $(0, 1/2)$.
The uniqueness of the root $r$ in $(0, 1/2)$ follows from the fact that the derivative of the  function \[Q'(r)=1+ \sum_{n=2}^{\infty} \left( 2-\frac{1}{n} \right)n r^{n-1}+ \sum_{n=2}^{\infty} \left( 2-\frac{1}{n} \right)^{p}pn r^{p n-1}>0.\]
 Thus,
\begin{equation*}
\frac{2 r+(1-r) \log(1-r)}{1-r}+\sum_{n=2}^{\infty} \left( 2-\frac{1}{n} \right)^{p} r^{pn} \le  (1 - \log{2}).
\end{equation*}
%Therefore, equation $(\ref{2.1.2})$ holds
 for $r \le r_p$ and $p \ge 1$, where $r_p$ is the solution of the equation $Q(r)=0$ in $(0, 1).$
To prove the sharpness, we assume
\begin{equation}\label{eq7}
    f(z)=\frac{2z}{(1+z)}- \log(1+z)= \sum_{n=1}^{\infty} (-1)^{n+1} \left(2-\frac{1}{n} \right) z^n.
\end{equation}
%From the inequality $(\ref{2.1.2})$, for $p \geq 1$
For  $z=r=r_p$, we have
\begin{align*}
|z|+ \sum_{n=2}^{\infty} |a_n z^n|& + \sum_{n=2}^{\infty} |a_n|^{p} |z|^{np}\\
&=|z|+ \sum_{n=2}^{\infty} \left| (-1)^{n+1} \left(2-\frac{1}{n} \right) z^n \right| + \sum_{n=2}^{\infty} {\left| (-1)^{n+1} \left(2-\frac{1}{n} \right) \right| }^p |z|^{np}\\
& =   |r|+ \sum_{n=2}^{\infty} \left| \left(2-\frac{1}{n} \right) r^n \right| + \sum_{n=2}^{\infty} {\left| \left(2-\frac{1}{n} \right) \right| }^p |r|^{np}\\
&= r+ \sum_{n=2}^{\infty}  \left(2-\frac{1}{n} \right) r^n  + \sum_{n=2}^{\infty} { \left(2-\frac{1}{n} \right)  }^p r^{np}\\
&=\frac{2 r+(1-r) \log(1-r)}{1-r}+\sum_{n=2}^{\infty} { \left(2-\frac{1}{n} \right)  }^p r^{np}\\
&=1-\log 2\\
& =\lim \limits_{r \rightarrow 1} \left(\frac{2r}{1+r}- \log(1+r)\right)\\
&= d(f(0), \partial{f(\mathbb{D})})
\end{align*}
Hence, the result is sharp.
\end{proof}

\begin{remark}
For different values of $p$ the radius converges to $0.215585$ upto $5$ decimal places as evident from the table given below
\end{remark}
\begin{table}[h!]
	\begin{center}
		\begin{tabular}{ |c|c|c| }
			\hline
			S. No. & p & $r_p$\\
			\hline
			1 & 2 & 0.213087\\
			\hline
			2 & 3 & 0.215411 \\
			\hline
			3 & 4 & 0.215573\\
			\hline
			4 & 5 & 0.215584\\
			\hline
			5 & 6& 0.215584\\
			\hline
			6 & 7& 0.215585\\
			\hline
			7& 8& 0.215585\\
			\hline
		\end{tabular}
		\caption{\textbf{Calculation of the  radius $r_p$ for various values of $p$.}}
		\label{Table:1}
	\end{center}
\end{table}
\noindent
In the  next two theorems \ref{th3} and \ref{th4}, we find the  Rogosinski's inequality and Rogosinski's radius for the class $\mathcal{C}_1$.
\begin{theorem}\label{th3}
Suppose that $f(z)=z + \sum_{n=2}^{\infty} a_n z^n \in \mathcal{C}_1$ and $N\geq2$ is an integer, then
\begin{equation} \label{2.1.3}
|f(z)|+ \sum_{n=N}^{\infty} |a_n z^n|  \le d(f(0), \partial{f(\mathbb{D})})
\end{equation}
for $|z| \le r_{1,N}$, where $r_{1,N}$ is the solution of
\begin{equation*}
\frac{2r}{1-r} +\log(1-r) + \sum_{n=N}^{\infty} \left( 2-\frac{1}{n} \right) r^{n}- (1 - \log{2})=0.
\end{equation*}
The result is sharp.
\end{theorem}
\begin{proof} Let $f\in  \mathcal{C}_1$. Using
  coefficients bound  ($\ref{coef1}$) and growth estimate ($\ref{growth1}$), for $|z| \le r$, we have
\begin{equation}\label{eq6}
|f(z)|+ \sum_{n=N}^{\infty} |a_n z^n|  \le
\frac{2r}{1-r} +\log(1-r) + \sum_{n=N}^{\infty} \left( 2-\frac{1}{n} \right) r^{n}.
\end{equation}
By \eqref{eq5} and \eqref{eq6}, the  inequality  holds
\begin{equation*}
|f(z)|+ \sum_{n=N}^{\infty} |a_n z^n|  \le  d(f(0), \partial{f(\mathbb{D})})
\end{equation*}
if and only if \begin{equation*}
\frac{2r}{1-r} +\log(1-r) + \sum_{n=N}^{\infty} \left( 2-\frac{1}{n} \right) r^{n} \le (1 - \log{2}),\quad |z|\leq r.
\end{equation*}
Let the function $R:[0, 1)\rightarrow \mathbb{R}$ be defined by
\begin{equation*}
R(r)=\frac{2r}{1-r} +\log(1-r) + \sum_{n=N}^{\infty} \left( 2-\frac{1}{n} \right) r^{n}- (1 - \log{2}).
\end{equation*}
Note that $R(0)=\log2-1<0$ and \[R\left(\frac{1}{2}\right)=2+\sum_{n=N}^{\infty} \left( 2-\frac{1}{n} \right) r^{n}- (1 - \log{2})>0.\]
Using the Intermediate Value Theorem, there exists a root $r\in(0, 1/2)$  of the equation $R(r)=0$. Also,
\begin{align*}
R'(r)&=\frac{2}{(1-r)^2}-\frac{1}{1-r}+\frac{2 r^{-1+N} (N+r-N r)}{(-1+r)^2}-\frac{r^{-1+N}}{1-r}\\
&=\frac{r+r^2+(2 N-1) r^N+(3-2 N) r^{1+N}}{(1-r)^2 r}\\
&\geq0.
\end{align*}
Therefore, the equation $R(r)=0$ has a unique root, say $r_{1,N}$. Thus, the  equation $(\ref{2.1.3})$ holds for $r \le r_{1,N}$.
The radius is sharp for the function $f$ defined in $\eqref{eq7}$.%
%  %  f(z)=\frac{2z}{(1+z)}- \log(1+z).
%\end{equation*}
%The power series expansion of $f(z)$ is the following
%\begin{equation*}
%    f(z)= \sum_{n=1}^{\infty} (-1)^{n+1} \left(2-\frac{1}{n} \right) z^n.
%\end{equation*}
%From the inequality $(\ref{2.1.3})$, we have
%\begin{equation*}
%|f(z)|+ \sum_{n=N}^{\infty} |a_n z^n|  =
%\left| \frac{2z}{1+z} -\log(1+z) \right| + \sum_{n=N}^{\infty} \left| (-1)^{n+1}\left( 2-\frac{1}{n} \right) z^{n} \right|
%\end{equation*}
%At $z=-r$
%\begin{equation*}
%|f(z)|+ \sum_{n=N}^{\infty} |a_n z^n|= \left| \frac{-2r}{1-r} -\log(1-r) \right| + \sum_{n=N}^{\infty} \left|\left( 2-\frac{1}{n} \right) (-r)^{n} \right|
%\end{equation*}
%\begin{equation*}
% =  \left| \frac{2r}{1-r} +\log(1-r) \right| + \sum_{n=N}^{\infty} \left|\left( 2-\frac{1}{n} \right) (-1)^n r^{n} \right|
%\end{equation*}
%\begin{equation} \label{2.1.3a}
%= \left| \frac{2r}{1-r} +\log(1-r) \right| + \sum_{n=N}^{\infty} \left|\left( 2-\frac{1}{n} \right) r^{n} \right|.
%\end{equation}
%Lastly,
%\begin{equation*}
%d(f(0), \partial{f(\Mathbb{D})}) = \liminf \limits_{z \rightarrow 1} \left(\frac{2z}{1+z}- \log(1+z)\right).
%\end{equation*}
%At $z=-r$
%\begin{equation}\label{2.1.3b}
% d(f(0), \partial{f(\Mathbb{D})}) = \liminf \limits_{r \rightarrow -1} \left(\frac{-2r}{1-r}- \log(1-r)\right)=1-\log 2.
%\end{equation}
%From equations $(\ref{2.1.3a})$ and $(\ref{2.1.3b})$, we get $f(z)$ satisfies
%\begin{equation*}
%\frac{2r}{1-r} +\log(1-r) + \sum_{n=N}^{\infty} \left( 2-\frac{1}{n} \right) r^{n}= (1 - \log{2}).
%\end{equation*}
%Hence, the result is sharp.
\end{proof}
\noindent
Next, we state the second result determining the sharp Bohr-Rogosinski radius for the class $\mathcal{C}_1$.
\begin{theorem}\label{th4}
Let $f(z)=z + \sum_{n=2}^{\infty} a_n z^n \in \mathcal{C}_1$ and $N\geq2$ be an integer. Then
\begin{equation} \label{2.1.4}
|f(z)|^2+ \sum_{n=N}^{\infty} |a_n z^n|  \le d(f(0), \partial{f(\mathbb{D})})
\end{equation}
for $|z| \le r_{2,N}$, where $r_{2,N}$ is the solution of
\begin{equation*}
\left(\frac{2r}{1-r} +\log(1-r) \right)^2 + \sum_{n=N}^{\infty} \left( 2-\frac{1}{n} \right) r^{n}- (1 - \log{2})=0.
\end{equation*}
The result is sharp.
\end{theorem}
\begin{proof}
%Following similar steps in previous theorems we begin by evaluating,
%\begin{equation*}
%d(f(0), \partial{f(\Mathbb{D})})= \liminf \limits_{|z| \rightarrow 1} |f(z)-f(0)|.
%\end{equation*}
%From growth estimate of $f(z) \in \mathcal{C}_1$ in equation ($\ref{growth1}$), we get
%\begin{equation*}
%d(f(0), \partial{f(\Mathbb{D})}) \ge \lim \limits_{r \rightarrow 1} \left(\frac{2r}{1+r}- \log(1+r)\right)
%\end{equation*}
%which gives $d(f(0), \partial{f(\Mathbb{D})}) \ge (1-\log{2})$ for $f \in \mathcal{C}_1$. Considering inequality in equation ($\ref{2.1.4}$), coefficient bound in equation ($\ref{coef1}$) and growth estimate in equation ($\ref{growth1}$)
 By Lemma \ref{lm1} and Lemma \ref{lm2}, for $|z| \le r$,
we have
\begin{equation*}
|f(z)|^2+ \sum_{n=N}^{\infty} |a_n z^n|  \le
 \left( \frac{2r}{1-r} +\log(1-r) \right)^2 + \sum_{n=N}^{\infty}  \left( 2-\frac{1}{n} \right) r^{n}.
\end{equation*}
To find the radius $r$ satisfying the inequality
\begin{equation*}
|f(z)|^2+ \sum_{n=N}^{\infty} |a_n z^n|  \le d(f(0), \partial{f(\mathbb{D})})\quad (|z|\leq r),
\end{equation*}
we define the function $S:[0, 1)\rightarrow\mathbb{R}$ such that
\begin{equation*}
S(r)=\left(\frac{2r}{1-r} +\log(1-r) \right)^2 + \sum_{n=N}^{\infty} \left( 2-\frac{1}{n} \right) r^{n}- (1 - \log{2}).
\end{equation*}
We claim that $S(r)\leq 0$ for $r \le r_{2,N}$, where $r_{2,N}$ is the solution of the equation
\begin{equation}\label{eq8}
\left(\frac{2r}{1-r} +\log(1-r) \right)^2 + \sum_{n=N}^{\infty} \left( 2-\frac{1}{n} \right) r^{n}- (1 - \log{2})=0.
\end{equation}
An easy calculation shows that $S(0)=\log2-1<0$ and
\[S\left(\frac{1}{2}\right)=4+\sum_{n=N}^{\infty} \left( 2-\frac{1}{n} \right) \frac{1}{2^{n}}- (1 - \log{2})
\geq4-(1-\log2)
=1+\log2
>0\]
 Clearly, by the Intermediate Value Theorem, there exists a root $r\in(0, 1/2)$  of the equation $S(r)=0$. Note that, for $r\in(0, 1/2)$,
 \begin{align*}S'(r)&=\left(\frac{2}{(1-r)^2}-\frac{1}{1-r}\right)^2+\frac{2 r^{N-1} (N+r-N r)}{(-1+r)^2}-\frac{r^{N-1}}{1-r}\\
 &=\left(\frac{2}{(1-r)^2}-\frac{1}{1-r}\right)^2+\frac{2N r^{N-1}(1-r) }{(1-r)^2}+\frac{2r^N}{(1-r)^2}-\frac{r^{N-1}}{1-r}\\
 &=\left(\frac{2}{(1-r)^2}-\frac{1}{1-r}\right)^2+\frac{2N r^{N-1}(1-r) }{(1-r)^2}+\frac{(3-r)r^{N-1}}{(1-r)^2}\\
 &\geq0
 \end{align*}
It shows that the function $S(r)$ is strictly increasing in $(0, 1/2)$. Thus, $S(r)\leq 0$ for $r \le r_{2,N}$, where $r_{2,N}$ is the unique root of the equation \eqref{eq8}.

To prove the sharpness, we assume
\begin{equation*}
    f(z)=\frac{2z}{(1+z)}- \log(1+z).
\end{equation*}

By \eqref{eq8}, at $z=-r=-r_{2,N}$, we have,
\begin{align*}
|f(z)|^2 + \sum_{n=N}^{\infty} |a_n z^n| & =
\left| \frac{2z}{1+z} -\log(1+z) \right|^2 + \sum_{n=N}^{\infty} \left| (-1)^{n+1}\left( 2-\frac{1}{n} \right) z^{n} \right|\\
&= \left| \frac{-2r}{1-r} -\log(1-r) \right|^2 + \sum_{n=N}^{\infty} \left|\left( 2-\frac{1}{n} \right) (-r)^{n} \right|\\
&= \left( \frac{2r}{1-r} +\log(1-r) \right)^2 + \sum_{n=N}^{\infty} \left( 2-\frac{1}{n} \right) r^{n}\\
&=(1 - \log{2})\\
&=\liminf \limits_{r \rightarrow 1} \left(\frac{2r}{1+r}- \log(1+r)\right)\\
&=d(f(0), \partial{f(\mathbb{D})}).
\end{align*}
Hence, the result is sharp.
\end{proof}

\section{certain Bohr radii for the class $\mathcal{C}_2$}
\noindent
We need the following lemmas for main results.
\begin{lemma} {\cite{Sil1972}}
If $f(z)= z + \sum_{n=2}^{\infty} a_n z^n \in \mathcal{C}_2$, then
\begin{equation}\label{coef2}
 |a_n| \le 1
\end{equation}
\begin{equation}\label{growth2}
\frac{r}{1+r}  \le |f(z)| \le \frac{r}{1-r}
\end{equation}
\begin{equation}\label{distortion2}
\frac{1}{(1+r)^2} \le |f^{'}(z)| \le \frac{1}{(1-r)^2}
\end{equation}
for all $|z| \le r$. Equality in all cases is obtained for $f(z)=z/(1-z)$.
\end{lemma}
\noindent
Using the above coefficient bound, growth and distortion results we obtain the following Bohr-type inequalities for the class $\mathcal{C}_2$.
\begin{theorem}\label{th1}
Suppose that $f(z)=z + \sum_{n=2}^{\infty} a_n z^n \in \mathcal{C}_2$, then
\begin{equation}\label{2.2.1}
|f(z)| + |f^{'}(z)||z| + \sum_{n=2}^{\infty} |a_n z^n| \le d( f(0), \partial{f(\mathbb{D})})
\end{equation}
for $|z| \le r_{21}\simeq 0.173417$, where $r_{21}$ is the solution of
\begin{equation*}
    1-6 r+r^2+2 r^3 =0.
\end{equation*}
The result is sharp.
\end{theorem}
\begin{proof}
We begin the proof by evaluating
\begin{equation*}
d(f(0), \partial{f(\mathbb{D})})= \liminf \limits_{|z| \rightarrow 1} |f(z)-f(0)|.
\end{equation*}
From growth estimate of $f(z) \in \mathcal{C}_2$ in equation ($\ref{growth2}$), we get
\begin{equation*}
d(f(0), \partial{f(\mathbb{D})}) \ge \lim \limits_{r \rightarrow 1} \left(\frac{r}{1+r}\right)
\end{equation*}
which gives
\begin{equation}\label{eq9}
d(f(0), \partial{f(\mathbb{D})}) \ge 1/2.
\end{equation}
\noindent
 Equations $(\ref{coef2})$, $(\ref{growth2})$ and $(\ref{distortion2})$ yield that for all $|z| \le r$, we have
\begin{equation}\label{eq10}
|f(z)| + |f^{'}(z)||z| + \sum_{n=2}^{\infty} |a_n| |z|^n
\le \frac{r}{1-r} + \frac{r}{(1-r)^2} + \sum_{n=2}^{\infty} r^{n}=\frac{r \left(2-r^2\right)}{(1-r)^2}.
\end{equation}
An easy calculation shows that
\begin{align}\label{eq11}
\frac{r \left(2-r^2\right)}{(1-r)^2}&\leq \frac{1}{2}\nonumber\\
\text{if and only if}\quad \frac{1-6 r+r^2+2 r^3}{2 (1-r)^2}&\geq0\nonumber\\
\text{if and only if} \quad1-6 r+r^2+2 r^3&\geq0.
\end{align}
By \eqref{eq10} and \eqref{eq11}, it follows that the inequality
\[|f(z)| + |f^{'}(z)||z| + \sum_{n=2}^{\infty} |a_n z^n| \le d( f(0), \partial{f(\mathbb{D})})\]
holds for all $|z|\leq r$, where $r=r_{21}\simeq0.173417\in(0, 1)$ is the root of the equation $1-6 r+r^2+2 r^3=0.$
\begin{figure}
  \centering
  % Requires \usepackage{graphicx}
  \includegraphics[width=10cm]{Thm32}\\
  \caption{The radius $r_{21}=0.173417$ is sharp}%\label{}
\end{figure}
\\
\noindent
To prove the sharpness, consider the function
\begin{equation*}
    f(z)=\frac{z}{(1-z)}=\sum_{n=1}^{\infty} z^n.
\end{equation*}
\\
\noindent
At $z=r_{21}$, we have
\begin{align*}
|f(z)| + |f^{'}(z)||z| + \sum_{n=2}^{\infty} |a_n| |z|^n &= \left|\frac{z}{(1-z)}\right|+ \left|\frac{z}{(1-z)^2}\right| +  \sum_{n=1}^{\infty}  |z|^n\\
&=\frac{r_{21}}{(1-r_{21})}+ \frac{r_{21}}{(1-r_{21})^2} +  \sum_{n=1}^{\infty}  r_{21}^n\\
&=\frac{r_{21} \left(2-r_{21}^2\right)}{(1-r_{21})^2}\\
&=\frac{1}{2}\\
&=d(f(0), \partial{f(\mathbb{D})}).
\end{align*}
Hence, the result is sharp.
\end{proof}
\noindent
Next, by considering powers of the coefficient $(|a_n|)$, we evaluate the corresponding sharp Bohr radii for the class $\mathcal{C}_2$.
\begin{theorem}
Suppose that $f(z)=z + \sum_{n=2}^{\infty} a_n z^n \in \mathcal{C}_2$, then
\begin{equation} \label{2.2.2}
|z|+ \sum_{n=2}^{\infty} |a_n z^n| + \sum_{n=2}^{\infty} |a_n|^{p} |z|^{np}  \le d(f(0), \partial{f(\mathbb{D})})
\end{equation}
for $|z| \le r^*_{p}$ and $p \ge 1$, where $r^*_{p}$ is the solution of
\begin{equation*}
1-3 r-r^p-2 r^{2 p}+3 r^{1+p}+2 r^{1+2 p}=0.
\end{equation*}
The result is sharp.
\end{theorem}
\begin{proof} Let $f\in\mathcal{C}_2.$
Now proceeding as  in the previous Theorem \ref{th1}, we get
\begin{equation*}
d(f(0), \partial{f(\mathbb{D})}) = \liminf \limits_{|z| \rightarrow 1} |f(z)-f(0)| \ge \lim \limits_{r \rightarrow 1} \left(\frac{r}{1+r}\right) = \frac{1}{2}
\end{equation*}
\\
\noindent
It follows from the equations $(\ref{coef2})$, $(\ref{growth2})$ and $(\ref{distortion2})$ that for all $|z| \le r$, $p \ge 1$,  the following inequality holds
\begin{equation*}
|z|+ \sum_{n=2}^{\infty} |a_n| |z|^n+ \sum_{n=2}^{\infty} |a_n|^{p} |z|^{np} \le r+ \sum_{n=2}^{\infty} r^{n}+ \sum_{n=2}^{\infty} r^{pn}=r+\frac{r^2}{1-r}+\frac{r^{2 p}}{1-r^p}.
\end{equation*}
Also, \[r+\frac{r^2}{1-r}+\frac{r^{2 p}}{1-r^p}=\frac{r+r^{2 p}-r^{1+p}-r^{1+2 p}}{(1-r) \left(1-r^p\right)}=\frac{1}{2}\]
if and only if $1-3 r-r^p-2 r^{2 p}+3 r^{1+p}+2 r^{1+2 p}=0$.
Let us assume that
\begin{equation*}
Q(r)=1-3 r-r^p-2 r^{2 p}+3 r^{1+p}+2 r^{1+2 p}.
\end{equation*}
Since, $Q(0)=1$ and
\noindent
\[Q(1/2)=2^{-1 - 2 p} (-2 + 2^p - 4^p)=-2^{-1 - 2 p} (2 - 2^p + 4^p)=-2^{-1 - 2 p} (2 + 2^p( 2^p-1)<0,\] by Intermediate Value Theorem, $Q(r)$ has a root in $(0, 1/2)$. To prove that the function $Q$ is decreasing  we need to find the maximum of the function
$j(p)=pr^p, (p\geq 1, r\in(0, 1/2))$. Clearly,
$j'(p)=r^p (1+p \log r).$ An easy calculation shows that  $j'(p)=0 $ if and only if $p= -{1}/{\log r}$. Also, $j''\left(-{1}/{\log r}\right)={\log r}/{e}<0$. Thus,
\begin{equation}\label{eq12}
pr^p\leq j(-{1}/{\log r})=-{1}/{(e \log r)}\leq -{1}/{(e \log (1/2))}={1}/{(e \log 2)}.\end{equation}
Using equation \eqref{eq12}, for $r\in(0, 1/2)$, we have
\begin{align*}
Q'(r)&=-3+3 r^p+2 r^{2 p}+p r^{p-1} \left(-1+3 r-4 r^p+4 r^{1+p}\right)\\
&<-3+\frac{3}{2}+\frac{1}{2}-p r^{p-1} +3pr^p -4p r^{2p-1}+4p r^{2p}\\
&=-1+pr^p\left(-\frac{1}{r}+3\right)+4pr^{2p}\left(-\frac{1}{r}+1\right)\\
&<-1+pr^p-4pr^{2p}\\
&<-1+\frac{1}{e \log 2}-4pr^{2p}\\
&<0.
\end{align*}
Thus, $Q(r)=0$ has the unique root $r\in (0, 1/2)$, say $r \le r^*_{p}$.
Therefore, equation $(\ref{2.2.2})$ holds for $r \le r^*_{p}$ and $p \ge 1$.
\\
\noindent
The result is sharp for the function
\begin{equation*}
    f(z)=\frac{z}{(1-z)}=\sum_{n=1}^{\infty} z^n.
\end{equation*}
%From the inequality $(\ref{2.2.2})$, for $p \geq 1 $
%\begin{equation*}
%|z|+ \sum_{n=2}^{\infty} |a_n z^n| + \sum_{n=2}^{\infty} |a_n|^{p} |z|^{np} = |z| + \sum_{n=2}^{\infty} |z^n| + \sum_{n=2}^{\infty} |z|^{np}.
%\end{equation*}
%At $z=-r$
%\begin{equation*} \label{2.2.2a}
%|z|+ \sum_{n=2}^{\infty} |a_n z^n| + \sum_{n=2}^{\infty} |a_n|^{p} |z|^{np} = |-r| + \sum_{n=2}^{\infty} |(-r)^n| + \sum_{n=2}^{\infty} |(-r)|^{np}.
%\end{equation*}
%Hence,
%\begin{equation}\label{2.2.2a}
%|z|+ \sum_{n=2}^{\infty} |a_n z^n| + \sum_{n=2}^{\infty} |a_n|^{p} |z|^{np} = |r| + \sum_{n=2}^{\infty} |r^n| + \sum_{n=2}^{\infty} |r|^{np}.
%\end{equation}
%Lastly, at $z=-r$
%\begin{equation} \label{2.2.2b}
%d(f(0), \partial{f(\mathbb{D})}) = \liminf \limits_{r \rightarrow -1} \left(\frac{-r}{1+r}\right)=\frac{1}{2}.
%\end{equation}
%From equations $(\ref{2.2.2a})$ and $(\ref{2.2.2b})$, we get
%\begin{equation*}
%r+ \sum_{n=2}^{\infty} r^{n}+ \sum_{n=2}^{\infty} r^{pn}= \frac{1}{2}.
%\end{equation*}
%Hence, the result is sharp.
\end{proof}
\begin{remark}
For different values of $p$ the radius converges to $0.333333$ upto $6$ decimal places as evident from the table given below
\end{remark}
\begin{table}[h!]
    \begin{center}
		\begin{tabular}{ |c|c|c|c|c|c|c|c|c|c| }
			\hline
			S. No. & 1 & 2&3&4&5&6&7\\
			\hline
			p & 2 & 3&4&5&6&7&8\\
			\hline
			$r^*_p$ & 0.327553 & 0.332707 & 0.333265 & 0.333326 & 0.333332 & 0.333333 &0.333333\\
			\hline
		\end{tabular}
        \caption{\textbf{Calculation of the  radius $r^*_p$ for various values of $p$.}}
		\label{Table:2}
    \end{center}
\end{table}

\noindent
We evaluate Bohr-Rogosinski radius for the class $\mathcal{C}_2$ in next two theorems.
\begin{theorem}
Suppose that $f(z)=z + \sum_{n=2}^{\infty} a_n z^n \in \mathcal{C}_2$ and  $N\geq2$ is an integer. Then
\begin{equation} \label{2.2.3}
|f(z)|+ \sum_{n=N}^{\infty} |a_n z^n|  \le d(f(0), \partial{f(\mathbb{D})})
\end{equation}
for $|z| \le r^*_{1,N}$, where $r^*_{1,N}$ is the solution of
\[ 3 r + 2 r^N-1=0.\]
The result is sharp.
\end{theorem}
\begin{proof}
Let $f\in  \mathcal{C}_2$. Using
  coefficients bound  ($\ref{coef2}$) and growth estimate ($\ref{growth2}$), for $|z| \le r$, we have
\begin{equation}\label{eq13}
|f(z)|+ \sum_{n=N}^{\infty} |a_n z^n|  \le
\frac{r}{1-r}  + \sum_{n=N}^{\infty} r^{n}.
\end{equation}
\noindent
By \eqref{eq9} and \eqref{eq13}, the  inequality  holds
\begin{equation*}
|f(z)|+ \sum_{n=N}^{\infty} |a_n z^n|  \le  d(f(0), \partial{f(\mathbb{D})})
\end{equation*}
if and only if
\begin{align*}
  \frac{r}{1-r}  + \sum_{n=N}^{\infty} r^{n}&=\frac{r+r^N}{1-r} \le   \frac{1}{2}
\end{align*}
\text{or equivalently} $ 3 r + 2 r^N-1=0.$
%As we have seen in the previous theorem, from growth estimate of $f(z) \in \mathcal{C}_2$ in equation ($\ref{growth2}$), we get
%\begin{equation*}
%d(f(0), \partial{f(\mathbb{D})}) = \liminf \limits_{|z| \rightarrow 1} |f(z)-f(0)| \ge \lim \limits_{r \rightarrow 1} \left(\frac{r}{1+r}\right) = \frac{1}{2}.
%\end{equation*}
%Therefore, $d(f(0), \partial{f(\mathbb{D})}) \geq 1/2$ for class $\mathcal{C}_2.$ Considering inequality in equation ($\ref{2.2.3}$), coefficient bound in equation ($\ref{coef2}$) and growth estimate in equation ($\ref{growth2}$) for $|z| \le r$ we get
%\begin{equation*}
%|f(z)|+ \sum_{n=N}^{\infty} |a_n z^n|  \le
% \frac{r}{1-r}  + \sum_{n=N}^{\infty} r^{n}.
%\end{equation*}
%Combining all the above inequalities we get equation ($\ref{2.2.3}$) holds if and only if
%\begin{equation*}
%  \frac{r}{1-r}  + \sum_{n=N}^{\infty} r^{n}  \le   \frac{1}{2}.
%\end{equation*}
Let us defined the function $R:[0, 1]\rightarrow 1$  such that
\begin{equation*}
R(r)= 3 r + 2 r^N-1.
\end{equation*}
Since $R(0)=-1$, $R(1)=4$ and $R'(r)=3+2 N r^{-1+N}>0$, using the  Intermediate Value Theorem, the equation $R(r)=0$ has the unique root $r\in(0, 1)$, say  $r^*_{1,N}$.
\\
\noindent
To prove the sharpness, we assume
\begin{equation*}
    f(z)=\frac{z}{(1-z)}=\sum_{n=1}^{\infty} z^n.
\end{equation*}
At $z=r=r^*_{1,N}$, we have
\begin{align*}
|f(z)|+ \sum_{n=N}^{\infty} |a_n z^n|&= \left| \frac{z}{1-z} \right|+ \sum_{n=N}^{\infty} |z|^n\\
&=\frac{r}{1-r}+ \sum_{n=N}^{\infty} r^n\\
&=\frac{r+r^N}{1-r}\\
&=\frac{1}{2}\\
&=d(f(0), \partial{f(\mathbb{D})}).
\end{align*}
Hence, the result is sharp.
\end{proof}
\begin{theorem}
Suppose that $f(z)=z + \sum_{n=2}^{\infty} a_n z^n \in \mathcal{C}_2$, then
\begin{equation} \label{2.2.4}
|f(z)|^2+ \sum_{n=N}^{\infty} |a_n z^n|  \le d(f(0), \partial{f(\mathbb{D})})
\end{equation}
for $|z| \le r^*_{2,N}$, where $r^*_{2,N}$ is the solution of
\begin{equation}\label{eq21}
1-2 r-r^2-2 r^N+2 r^{1+N}=0.
\end{equation}
The result is sharp.
\end{theorem}
\begin{proof}For $|z| \le r$,
the coefficient bounds  ($\ref{coef2}$) and growth estimate  ($\ref{growth2}$) yield
\begin{align*}
|f(z)|^2+ \sum_{n=N}^{\infty} |a_n z^n|&  \le \left( \frac{r}{1-r} \right)^2 + \sum_{n=N}^{\infty} r^{n}\nonumber\\
&=\frac{r^2+r^N-r^{1+N}}{(1-r)^2}.
\end{align*}
Also, by \eqref{eq9}
\[
d(f(0), \partial{f(\mathbb{D})}) \ge 1/2.
\]
Thus, ($\ref{2.2.4}$) holds if and only if
\begin{equation*}
 \frac{r^2+r^N-r^{1+N}}{(1-r)^2}   \le   \frac{1}{2}
\end{equation*}
equivalently
\[-1+2 r+r^2+2 r^N-2 r^{1+N}\leq0.\]
Let us define a function
\begin{equation*}
S(r)=-1+2 r+r^2+2 r^N-2 r^{1+N}\quad (r\in[0,1 )).
\end{equation*}
It is easy to see that
$S(0)=-1<0$,  $S(1/2)=1/4+1/2^N>0$ and
\begin{align*}
S'(r)&= \frac{2 \left(r+r^2+N r^N-r^{1+N}-N r^{1+N}\right)}{r}\\
&=\frac{2 \left(r+r^2+ r^N(N-(1+N) r)\right)}{r}\\
&>\frac{2 \left(r+r^2- r^N\right)}{r}\\
&=\frac{2 \left(r(1-r^{N-1})+r^2\right)}{r}\\
&>0.
\end{align*}
In view of the above inequalities and the Intermediate Value Theorem, there exists the unique  root $r^*_{2,N}\in(0, 1/2)$ of the equation \eqref{eq21} which also satisfy the inequality $(\ref{2.2.4})$ for $r \le r^*_{2,N}$.
\\
\noindent
The result is sharp for the function
\begin{equation*}
f(z)=\frac{z}{(1-z)}=\sum_{n=1}^{\infty} z^n.
\end{equation*}
\end{proof}

\section{certain Bohr radii for the class $\mathcal{C}_3$}
\noindent
We have sharp coefficient bounds and distortion theorems for the class $\mathcal{C}_3$.
\begin{lemma} {\cite{Sil1972}}\label{lm3}
If $f(z)= z + \sum_{n=2}^{\infty} a_n z^n \in \mathcal{C}_3$, then
\begin{equation}\label{coef3}
 |a_n| \le \frac{2}{3} + \frac{1}{3n^2}  .
\end{equation}
This result is sharp, with equality for
\begin{equation*}
    f(z)= \frac{2z}{3(1-z)} - \frac{1}{3} \int_{0}^{z} \frac{\log(1-\zeta)}{\zeta} d\zeta.
\end{equation*}
\end{lemma}
\begin{lemma} {\cite{Sil1972}}\label{lm4}
If $f \in \mathcal{C}_3$ then
\begin{equation}
\frac{2r}{3(1+r)}+\frac{1}{3} \int_{0}^{r} \frac{\log(1+t)}{t} dt
\le |f(z)|
\le \frac{2r}{3(1-r)}-\frac{1}{3} \int_{0}^{r} \frac{\log(1-t)}{t} dt
\label{growth3}
\end{equation}
and
\begin{equation}
\frac{2}{3(1+r)^2}+ \frac{\log(1+r)}{3r}
\le |f^{'}(z)| \le
\frac{2}{3(1-r)^2}- \frac{\log(1-r)}{3r}
\label{distortion3}
\end{equation}
for all $(0 < |z| \le r)$ and equality holds in all cases for the extremal function
\begin{equation*}
    f(z)= \frac{2z}{3(1-z)} - \frac{1}{3} \int_{0}^{z} \frac{\log(1-\zeta)}{\zeta} d\zeta.
\end{equation*}
\end{lemma}
\noindent
Using the above coefficient bound, growth and distortion results we obtain the following Bohr-type inequalities for the class $\mathcal{C}_3$.
\begin{theorem}\label{th2}
Suppose that $f(z)=z + \sum_{n=2}^{\infty} a_n z^n \in \mathcal{C}_3$, then
\begin{equation}\label{2.3.1}
|f(z)| + |f^{'}(z)||z| + \sum_{n=2}^{\infty} |a_n z^n| \le d( f(0), \partial{f(\mathbb{D})})
\end{equation}
for $|z| \le r_{31}$, where $r_{31}$ is the solution of
\begin{equation*}
\frac{2 (2-r^2) r}{3 (1-r)^2}-\frac{1}{3} \int_{0}^{r} \frac{\log(1-t)}{t} dt-\frac{1}{3} \log(1-r)+ \sum_{n=2}^{\infty} \frac{r^n}{3 n^2}-\left(\frac{1}{3}+\frac{\pi ^2}{36}\right)=0.
\end{equation*}
The radius is sharp.
\end{theorem}
\begin{proof}
We begin the proof by evaluating
\begin{equation*}
d(f(0), \partial{f(\mathbb{D})})= \liminf \limits_{|z| \rightarrow 1} |f(z)-f(0)|.
\end{equation*}
From equation ($\ref{growth3}$), we get
\begin{equation*}
d(f(0), \partial{f(\mathbb{D})}) \ge \lim \limits_{r \rightarrow 1} \left(\frac{2r}{3(1+r)}+\frac{1}{3} \int_{0}^{r} \frac{\log(1+t)}{t} dt \right)
\end{equation*}
which gives
\begin{equation}\label{eq16}
d(f(0), \partial{f(\mathbb{D})}) \ge \frac{1}{3} +\frac{1}{3} \int_{0}^{1} \frac{\log(1+t)}{t} dt=\frac{1}{3}+\frac{\pi ^2}{36}.
\end{equation}
From  equations $(\ref{coef3})$, $(\ref{growth3})$ and $(\ref{distortion3})$, for all $|z| \le r$, we have
\begin{align}\label{eq17}
|f(z)|& + |f^{'}(z)||z| + \sum_{n=2}^{\infty} |a_n| |z|^n\\
&\le \frac{2r}{3(1-r)}-\frac{1}{3} \int_{0}^{r} \frac{\log(1-t)}{t} dt+ \frac{2r}{3(1-r)^2}-\frac{1}{3} \log(1-r)+ \sum_{n=2}^{\infty} \left( \frac{2}{3} + \frac{1}{3 n^2} \right)  r^n\nonumber\\
&=\frac{2 (2-r^2) r}{3 (1-r)^2}-\frac{1}{3} \int_{0}^{r} \frac{\log(1-t)}{t} dt-\frac{1}{3} \log(1-r)+ \sum_{n=2}^{\infty} \frac{r^n}{3 n^2}.
\end{align}
In order to find the radius $r_{31}\in (0, 1)$ satisfying the inequality $(\ref{2.3.1})$, we define the function $P:[0, 1)\rightarrow\mathbb{R}$ by

\begin{equation*}
P(r)=\frac{2 (2-r^2) r}{3 (1-r)^2}-\frac{1}{3} \int_{0}^{r} \frac{\log(1-t)}{t} dt-\frac{1}{3} \log(1-r)+ \sum_{n=2}^{\infty} \frac{r^n}{3 n^2}-\left(\frac{1}{3}+\frac{\pi ^2}{36}\right)
\end{equation*}
An easy calculation shows that
\begin{equation}\label{eq14}
P(0)=-\left(\frac{1}{3}+\frac{\pi ^2}{36}\right)<0
\end{equation}
and
\begin{align}\label{eq15}
P\left(\frac{1}{2}\right)&=2-\frac{\pi ^2}{36}+\frac{\log2}{3}+\frac{1}{36} \left(-6+\pi ^2-6 \log 2^2\right)+\frac{1}{36} \left(\pi ^2-6 \log 2^2\right)\nonumber\\
&=\frac{1}{36} \left(66+\pi ^2-12 \log 2^2+\log4096\right)\nonumber\\
&\simeq2.17839>0.
\end{align}
Equations \eqref{eq14} and \eqref{eq15} together with Intermediate Value Theorem yield that there exists a root $r\in(0, 1/2)$ of the equation $P(r)=0.$
Also,
\begin{align*}
P'(r)&=\frac{2 \left(2+2 r-3 r^2+r^3\right)}{3 (1-r)^3}-\frac{\log(1-r)}{3r}+\frac{1}{3 (1-r)}-\frac{1}{3} \left(1+\frac{\log(1-r)}{r}\right)\\
&=\frac{1}{3} \left(\frac{4-r (-5+(8-3 r) r)}{(1-r)^3}-\frac{2 \log(1-r)}{r}\right)\\
&=\frac{r \left(4+5 r-8 r^2+3 r^3\right)-2 (1-r)^3 \log(1-r)}{3 (1-r)^3 r}\\
&=\frac{r \left(4+(1-r) r (5-3 r)\right)-2 (1-r)^3 \log(1-r)}{3 (1-r)^3 r}\\
&>0.
\end{align*}
Thus, there exists the  unique  root $r\in(0, 1)$ of the equation $P(r)=0$, say $r_{31}$. By equations \eqref{eq16} and \eqref{eq17}, for $|z|\leq r_{31}$, we have
\[|f(z)| + |f^{'}(z)||z| + \sum_{n=2}^{\infty} |a_n z^n| \le d( f(0), \partial{f(\mathbb{D})}).\]
To prove the sharpness, consider the function
\begin{equation*}
    f(z)= \frac{2z}{3(1-z)} - \frac{1}{3} \int_{0}^{z} \frac{\log(1-\zeta)}{\zeta} d\zeta.
\end{equation*}
Note that
\[-\frac{1}{3}\int_0^r \frac{\log(1-t)}{t} \, dt=\frac{r}{3}+\frac{r^2}{12}+\frac{r^3}{27}+\frac{r^4}{48}+\frac{r^5}{75}+\frac{r^6}{108}+\cdots\]
and \[-(1/3) \log(1 - r)=\frac{r}{3}+\frac{r^2}{6}+\frac{r^3}{9}+\frac{r^4}{12}+\frac{r^5}{15}+\frac{r^6}{18}+\cdots\]
So, at  $z=r=r_{31}$, we have
\begin{align*}
|f(z)|& + |f^{'}(z)||z| + \sum_{n=2}^{\infty} |a_n| |z|^n\\
&= \left|\frac{2r}{3(1-r)} - \frac{1}{3} \int_{0}^{r} \frac{\log(1-\zeta)}{\zeta} d\zeta \right| + |r| \left| \frac{2}{3(1-r)^2}-\frac{\log(1-r)}{3r} \right|\\
  &+   \sum_{n=2}^{\infty} \left| \left( \frac{2}{3}+ \frac{n^2}{3}\right) \right| |r|^n\\
&= \frac{2r}{3(1-r)} - \frac{1}{3} \int_{0}^{r} \frac{\log(1-\zeta)}{\zeta} d\zeta  +   \frac{2r}{3(1-r)^2}-\frac{\log(1-r)}{3r}  +   \sum_{n=2}^{\infty}  \left( \frac{2}{3}+ \frac{n^2}{3}\right) r^n\\
&=\frac{2 (2-r^2) r}{3 (1-r)^2}-\frac{1}{3} \int_{0}^{r} \frac{\log(1-t)}{t} dt-\frac{1}{3} \log(1-r)+ \sum_{n=2}^{\infty} \frac{r^n}{3 n^2}\\
&=\frac{1}{3}+\frac{\pi ^2}{36}\\
&=\frac{1}{3} +\frac{1}{3} \int_{0}^{1} \frac{\log(1+t)}{t} dt\\
&=d(f(0), \partial{f(\mathbb{D})}).
\end{align*}
Hence, the result is sharp.
\end{proof}
\noindent
Next we prove the following Bohr radius associated with powers of the coefficient $(|a_n|)$ for the class $\mathcal{C}_3$.
\begin{theorem}
Suppose that $f(z)=z + \sum_{n=2}^{\infty} a_n z^n \in \mathcal{C}_3$, then
\begin{equation} \label{2.3.2}
|z|+ \sum_{n=2}^{\infty} |a_n z^n| + \sum_{n=2}^{\infty} |a_n|^{p} |z|^{np}  \le d(f(0), \partial{f(\mathbb{D})})
\end{equation}
for $|z| \le R_{p}$ and $p \ge 1$, where $R_{p}\in(0, 1/2)$ is the root of the equation
\begin{equation}\label{eq20}
\frac{(3-r) r}{3 (1-r)}+\sum_{n=2}^{\infty} \frac{r^n}{3 n^2}  + \sum_{n=2}^{\infty} \left( \frac{2}{3} + \frac{1}{3 n^2} \right)^p  r^{np}- \frac{1}{36} \left(12+\pi ^2\right)=0.
\end{equation}
The result is sharp.
\end{theorem}
\begin{proof}
Proceeding as in Theorem \ref{th2}, we have
\begin{equation}\label{eq18}
 d(f(0), \partial{f(\mathbb{D})}) \ge  \frac{1}{36} \left(12+\pi ^2\right).
\end{equation}
It follows from  equations $(\ref{coef3})$ that for $|z| \le r$, we have
\begin{align}\label{eq19}
|z|+ \sum_{n=2}^{\infty} |a_n| |z|^n+ \sum_{n=2}^{\infty} |a_n|^{p} |z|^{np}
&\le r+ \sum_{n=2}^{\infty} \left( \frac{2}{3} + \frac{1}{3 n^2} \right)  r^n + \sum_{n=2}^{\infty} \left( \frac{2}{3} + \frac{1}{3 n^2} \right)^p  r^{np}\nonumber\\
&=\frac{(3-r) r}{3 (1-r)}+\sum_{n=2}^{\infty} \frac{r^n}{3 n^2}  + \sum_{n=2}^{\infty} \left( \frac{2}{3} + \frac{1}{3 n^2} \right)^p  r^{np}
\end{align}
In view of the equations \eqref{eq18} and \eqref{eq19}, the inequality ($\ref{2.3.2}$) holds if and only if
\begin{equation*}
\frac{(3-r) r}{3 (1-r)}+\sum_{n=2}^{\infty} \frac{r^n}{3 n^2}  + \sum_{n=2}^{\infty} \left( \frac{2}{3} + \frac{1}{3 n^2} \right)^p  r^{np}\le  \frac{1}{36} \left(12+\pi ^2\right).
\end{equation*}
Now, we define a function $Q:[0, 1)\rightarrow\mathbb{R}$ such that
\begin{equation*}
Q(r)=\frac{(3-r) r}{3 (1-r)}+\sum_{n=2}^{\infty} \frac{r^n}{3 n^2}  + \sum_{n=2}^{\infty} \left( \frac{2}{3} + \frac{1}{3 n^2} \right)^p  r^{np}-\frac{1}{36} \left(12+\pi ^2\right).
\end{equation*}
An easy calculation shows that
\[Q(0)=-\frac{1}{36} \left(12+\pi ^2\right)<0\]
and
\begin{align*}
Q\left(\frac{1}{2}\right)&=\frac{1}{3}-\frac{(\log2)^2}{6}+\sum _{n=2}^{\infty } \left(\frac{2}{3}+\frac{2^{-n}}{3 n^2}\right)^p\frac{1}{2^{n p}}\nonumber\\
&=\frac{1}{6}\left(2-(\log2)^2\right)+\sum _{n=2}^{\infty } \left(\frac{2}{3}+\frac{2^{-n}}{3 n^2}\right)^p\frac{1}{2^{n p}}\nonumber\\
&>0.
\end{align*}
Since $Q(0).Q(1/2)<0$, by Intermediate Value Theorem, there exists a root $r\in(0, 1/2)$ satisfying the equation $Q(r)=0$. Moreover,
\begin{align*}
Q'(r)&=\frac{3-2 r+r^2}{3 (1-r)^2}-\frac{r+\log(1-r)}{3 r}+\sum _{n=2}^{\infty } \left(\frac{2}{3}+\frac{1}{3 n^2}\right)^p {np} r^{{np}-1}>0.
\end{align*}
Thus, there exists the unique root, say $R_{p}\in(0, 1/2)$ satisfying the equation \eqref{eq20}.
In order to prove the sharpness, consider the function
\begin{equation*}
    f(z)= \frac{2z}{3(1-z)} - \frac{1}{3} \int_{0}^{z} \frac{\log(1-\zeta)}{\zeta} d\zeta.
\end{equation*}
At $z=r=R_{p}$, we have
\begin{align*}
|z|+ \sum_{n=2}^{\infty} |a_n z^n| + \sum_{n=2}^{\infty} |a_n|^{p} |z|^{np}
&=|z|+ \sum_{n=2}^{\infty} \left|\left( \frac{2}{3} + \frac{1}{3 n^2} \right) z^n\right| + \sum_{n=2}^{\infty} \left|\left( \frac{2}{3} + \frac{1}{3 n^2} \right)\right|^{p} |z|^{np}\\
&=r+ \sum_{n=2}^{\infty} \left( \frac{2}{3} + \frac{1}{3 n^2} \right) r^n + \sum_{n=2}^{\infty} \left( \frac{2}{3} + \frac{1}{3 n^2} \right)^{p} r^{np}\\
&= \frac{1}{3} + \frac{1}{3} \int_{0}^{1} \frac{\log(1-t)}{t} dt\\
&=d(f(0), \partial{f(\mathbb{D})}).
\end{align*}
Hence, the result is sharp.
\end{proof}
\noindent
We find Bohr-Rogosinski radius for the class $\mathcal{C}_3$ in next two theorems.
\begin{theorem}
Suppose that $f(z)=z + \sum_{n=2}^{\infty} a_n z^n \in \mathcal{C}_3$, then
\begin{equation} \label{2.3.3}
|f(z)|+ \sum_{n=N}^{\infty} |a_n z^n|  \le d(f(0), \partial{f(\mathbb{D})})
\end{equation}
for $|z| \le R_{1,N}$, where $R_{1,N}$ is the solution of
\begin{equation*}
 \sum_{n=N}^{\infty}   \frac{r^{n}}{3n^2} -\frac{1}{3} \int_{0}^{r} \frac{\log(1-t)}{t} dt- \frac{12+\pi ^2-36 r-\pi ^2 r-24 r^N}{36 (1-r)}=0.
\end{equation*}
The result is sharp.
\end{theorem}
\begin{proof}
%Following similar steps in previous theorems, for $f \in \mathcal{C}_3$ we have
%\begin{equation*}
%d(f(0), \partial{f(\mathbb{D})}) \ge   \frac{1}{3} +\frac{1}{3} \int_{0}^{1} \frac{\log(1+t)}{t} dt.
%\end{equation*}
Using Lemma \ref{lm3} and growth estimate  ($\ref{growth3}$), for $|z| \le r$, we get
\begin{equation*}
|f(z)|+ \sum_{n=N}^{\infty} |a_n z^n|  \le \frac{2r}{3(1-r)}-\frac{1}{3} \int_{0}^{r} \frac{\log(1-t)}{t} dt+ \sum_{n=N}^{\infty} \left( \frac{2}{3} + \frac{1}{3n^2} \right) r^{n}.
\end{equation*}
It follows from  the above inequality that the  equation ($\ref{2.3.3}$) holds if and only if
\begin{equation*}
\frac{2r}{3(1-r)}-\frac{1}{3} \int_{0}^{r} \frac{\log(1-t)}{t} dt+ \sum_{n=N}^{\infty} \left( \frac{2}{3} + \frac{1}{3n^2} \right) r^{n} \le  \frac{1}{3} +\frac{1}{3} \int_{0}^{1} \frac{\log(1+t)}{t} dt.
\end{equation*}
In order to find the radius $R_{1,N}\in (0, 1)$ satisfying the inequality $(\ref{2.3.3})$, we define the function $R:[0, 1)\rightarrow\mathbb{R}$ by
\begin{equation*}
R(r)=\sum_{n=N}^{\infty}   \frac{r^{n}}{3n^2} -\frac{1}{3} \int_{0}^{r} \frac{\log(1-t)}{t} dt- \frac{12+\pi ^2-36 r-\pi ^2 r-24 r^N}{36 (1-r)}.
\end{equation*}
It is easy to see that
$R(0)=-1/3-\pi^2/36<0$ and
\begin{align*}
R\left(\frac{1}{2}\right)&=\sum_{n=N}^{\infty}\frac{1}{3n^2}\left(\frac{1}{2^{n}}\right)+\frac{1}{36} \left(\pi ^2-6 (\log2)^2\right)+\frac{2}{3} \left(1+2^{1-N}\right)-\frac{1}{36} \left(12+\pi ^2\right)\\
&=\sum_{n=N}^{\infty}\frac{1}{3n^2}\left(\frac{1}{2^{n}}\right)+\frac{1}{6} \left(2+2^{3-N}-(\log 2)^2\right)\\
&>0.
\end{align*}
In view of the above and using Intermediate Value Theorem it follows that there exists a root $\varepsilon\in(0, 1/2)$ such that $R(\varepsilon)=0.$
Since \
\[
R'(r)=\frac{2}{3 (1-r)^2}-\frac{\log(1-r)}{3 r}+\frac{2 r^{N-1} (N+r-N r)}{3 (1-r)^2}+\sum_{n=N}^{\infty}   \frac{r^{n-1}}{3n}>0,
\] there exists the unique $ R_{1,N}\in (0, 1/2)$ such that $R(R_{1,N})=0.$
\\
\noindent
Thus, equation $(\ref{2.3.3})$ holds for $r \le R_{1,N}$, where $R_{1,N}$ is the solution of the equation
\begin{equation*}
\frac{2r}{3(1-r)}-\frac{1}{3} \int_{0}^{r} \frac{\log(1-t)}{t} dt+ \sum_{n=N}^{\infty} \left( \frac{2}{3} + \frac{1}{3n^2} \right) r^{n} -\frac{1}{3}- \frac{1}{3} \int_{0}^{1} \frac{\log(1+t)}{t} dt=0.
\end{equation*}
To prove the sharpness, we assume
\begin{equation*}
   f(z)= \frac{2z}{3(1-z)} - \frac{1}{3} \int_{0}^{z} \frac{\log(1-\zeta)}{\zeta} d\zeta.
\end{equation*}
At $z=r=R_{1,N}$, we have
\begin{align*}%\label{2.3.3a}
|f(z)|+ \sum_{n=N}^{\infty} |a_n z^n|&=  \left|\frac{2z}{3(1-z)} - \frac{1}{3} \int_{0}^{z} \frac{\log(1-t)}{t} dt \right|+\sum_{n=N}^{\infty} \left|\left( \frac{2}{3} + \frac{1}{3 n^2} \right)\right||z|^{n}\\
 &=\frac{2r}{3(1-r)} - \frac{1}{3} \int_{0}^{r} \frac{\log(1-t)}{t} dt +\sum_{n=N}^{\infty} \left( \frac{2}{3} + \frac{1}{3 n^2} \right)r^{n}\\
 &=\frac{1}{3} + \frac{1}{3} \int_{0}^{1} \frac{\log(1-t)}{t} dt\\
 &=d(f(0), \partial{f(\mathbb{D})}).
\end{align*}
Hence, the result is sharp.
\end{proof}
\begin{theorem}
Suppose that $f(z)=z + \sum_{n=2}^{\infty} a_n z^n \in \mathcal{C}_3$, then
\begin{equation} \label{2.3.4}
|f(z)|^2+ \sum_{n=N}^{\infty} |a_n z^n|  \le d(f(0), \partial{f(\mathbb{D})})
\end{equation}
for $|z| \le R_{2,N}$, where $R_{2,N}$ is the solution of
\begin{equation*}
\left(
\frac{2r}{3(1-r)}-\frac{1}{3} \int_{0}^{r} \frac{\log(1-t)}{t} dt
\right)^2
+ \sum_{n=N}^{\infty}   \frac{r^{n}}{3n^2}  +\frac{ 24 r^N+(12 +\pi ^2)r-12-\pi ^2  }{36 (1-r)}=0.
\end{equation*}
The result is sharp.
\end{theorem}
\begin{proof}
It follows from Lemma \ref{lm3} and Lemma \ref{lm4} that
\begin{align*}
|f(z)|^2+ \sum_{n=N}^{\infty} |a_n z^n| & \le \left( \frac{2r}{3(1-r)}-\frac{1}{3} \int_{0}^{r} \frac{\log(1-t)}{t} dt \right)^2 + \sum_{n=N}^{\infty} \left( \frac{2}{3} + \frac{1}{3n^2} \right) r^{n}\\
&=\left( \frac{2r}{3(1-r)}-\frac{1}{3} \int_{0}^{r} \frac{\log(1-t)}{t} dt \right)^2 + \sum_{n=N}^{\infty} \frac{r^{n}}{3n^2}+\frac{2 r^N}{3(1-r)}
\end{align*}
for $|z| \le r$.
\\
\noindent
Define a function $S:[0, 1)\rightarrow\mathbb{R}$ by
\begin{equation*}
S(r)=\left( \frac{2r}{3(1-r)}-\frac{1}{3} \int_{0}^{r} \frac{\log(1-t)}{t} dt \right)^2+ \sum_{n=N}^{\infty} \left( \frac{2}{3} + \frac{1}{3n^2} \right) r^{n} -\frac{1}{3}- \frac{1}{3} \int_{0}^{1} \frac{\log(1+t)}{t} dt.
\end{equation*}
Note that
$S(0)=-{1}/{3}-{\pi ^2}/{36}$.
Also, an easy calculation yields
\begin{align*}
S(1/2)&=\left( \frac{2}{3}-\frac{1}{3} \int_{0}^{1/2} \frac{\log(1-t)}{t} dt \right)^2+ \sum_{n=N}^{\infty}  \frac{1}{3n^22^n}  +\frac{2^{2-N}}{3}-\frac{1}{36} \left(12+\pi ^2\right)\\
&=\frac{\left(24+\pi ^2-6 (\log2)^2\right)^2}{1296}+ \sum_{n=N}^{\infty}  \frac{1}{3n^22^n}  +\frac{2^{2-N}}{3}-\frac{1}{36} \left(12+\pi ^2\right)\\
&=\frac{\left(24+\pi ^2-6 (\log2)^2\right)^2}{1296}-\frac{1}{36} \left(12+\pi ^2\right)+ \sum_{n=N}^{\infty}  \frac{1}{3n^22^n}  +\frac{2^{2-N}}{3}\\
&=\frac{144+12 \pi ^2+\pi ^4-288 (\log2)^2-12 \pi ^2 (\log2)^2+36 (\log2)^4}{1296}+ \sum_{n=N}^{\infty}  \frac{1}{3n^22^n}  +\frac{2^{2-N}}{3}\\
&>0.
\end{align*}
In view of the above, $S(0)S(1/2)<0$, and by the  Intermediate Value Theorem there exists a root $\zeta\in(0, 1/2)$ such that $S(r)\leq 0$ for $r\leq \zeta$. Furthermore, Since
\begin{align*}
S'(r)&=2\left( \frac{2r}{3(1-r)}-\frac{1}{3} \int_{0}^{r} \frac{\log(1-t)}{t} dt \right)\left(\frac{2}{3 (1-r)^2}-\frac{\log(1-r)}{3 r}\right)\\
&+\frac{2 r^{-1+N} (N+r-N r)}{3 (1-r)^2}+\sum_{n=N}^{\infty}   \frac{r^{n-1}}{3n}\\
&>0,
\end{align*}
there exists the unique $ R_{2,N}\in(0, 1/2)$ such that the inequality $(\ref{2.3.4})$ holds for $r \le R_{2,N}$, and satisfying the equation $S( R_{2,N})=0.$
\\
\noindent
The result is sharp for the function
\begin{equation*}
   f(z)= \frac{2z}{3(1-z)} - \frac{1}{3} \int_{0}^{z} \frac{\log(1-\zeta)}{\zeta} d\zeta.
\end{equation*}
\end{proof}

\end{section}

%\end{section}

\noindent
\textbf{Data availability statement} Data sharing not applicable to this article as no
datasets were generated or analysed during the current study.
%

%are not compulsory. Where included they should be brief. Grant or contribution numbers may be acknowledged.
%
%Please refer to Journal-level guidance for any specific requirements.

\section*{Declarations}

%
%Some journals require declarations to be submitted in a standardised format. Please check the Instructions for Authors of the journal to which you are submitting to see if you need to complete this section. If yes, your manuscript must contain the following sections under the heading `Declarations':
\noindent
\textbf{Conflict of interest}  The authors declare they have no conflict of interest regarding the publication of this paper.

\noindent
\textbf{Funding}
No funding from any source.\\


































































\bibliographystyle{abbrv}
\bibliography{NJ-bibliography}

%\bibitem{bohr}H. Bohr, A Theorem Concerning Power Series, Proc. London Math. Soc. (2) {\bf 13} (1914), 1--5.
%
%\bibitem{goodman1}  A.~W. Goodman, {\it Univalent functions. Vol. I}, Mariner, Tampa, FL, 1983.
%
%\bibitem{goodman2}  A.~W. Goodman, {\it Univalent functions. Vol. II}, Mariner, Tampa, FL, 1983.
%
%\bibitem{silverman} H. Silverman and D. Telage, Extreme points of subclasses of close-to-convex functions, Proc. Amer. Math. Soc. {\bf 74} (1979), no.~1, 59--65.
%
%\bibitem{rogosinski} ~W. Rogosinski, \"Uber Bildschranken bei Potenzreihen und ihren Abschnitten, Math. Z. {\bf 17} (1923), no.~1, 260--276.
%
%\bibitem{ponnu} Kayumov, I.R. and Ponnusamy, S., 2017. Bohr--Rogosinski radius for analytic functions, preprint arXiv:1708.05585.

%\end{thebibliography}
\end{document}
